%% Part V, school Report #5. Grading: 20% on-going, 20% defence. Budget: 7 pages.
\chapter{System Design and Implementation}\label{ch:design}

\section{AI Model Integration}\label{sec:integration}
OpenTrading does not contain a model of its own.
It is a harness: the program around a large language model that decides what the model sees, which actions it may take and when a task is finished.
The model is reached through the interface of a model provider, and the market data is reached through the Model Context Protocol (MCP), an open standard for connecting models to tools and data.
\todo{Cite the MCP specification after opening it.}
\Cref{tab:components} lists the main components and the role of each.
\begin{table}[htb]
  \centering
  \caption{Main components of the system and their roles}\label{tab:components}
  \begin{tabularx}{\textwidth}{@{}l X@{}}
    \toprule
    {Component} & {Role}\\
    \midrule
    Language model & Reads the conversation and the tool results, and decides the next step or the final answer\\
    OpenTrading harness & Runs the agent loop and offers 23 built-in tools, 43 skills and 17 specialised sub-agents\\
    MCP client & Part of the harness; connects to one or more MCP servers and turns their tools into tools that the model can call\\
    opentrading-mcp & Data server built by the group; offers 175 tools and 6 analysis workflows over the market data\\
    Data providers & External sources of prices, financial statements, global markets and news (\cref{sec:market-data})\\
    \bottomrule
  \end{tabularx}
  \fignote{Source: OpenTrading \texttt{7319264} (version 0.5.1) and opentrading-mcp \texttt{de956f8}.}
\end{table}
%% Evidence: E10, E12, E21.

\subsection{From Server Tools to Model Tools}\label{sec:tool-exposure}
Each tool of the server is published with a name, a description, the schema of its inputs and a short search hint.
The MCP client gives the model the tools of every connected server next to the harness's own tools, with the server's name as a prefix so that two servers cannot clash.
A full list of 175 tool schemas would take a large part of the model's context in every request.
When the number of tools passes a threshold, the harness therefore sends only their names, and the model loads the tools it needs through a search tool, as described in \cref{sec:algorithm}.
The workflows of the server appear to the user as commands, so that a prepared analysis, such as a risk check of a portfolio, can be started in one line.

\subsection{A Shared Answer Contract}\label{sec:answer-contract}
A tool call can fail for reasons outside the model's control: a provider is down, a market is closed, or a company has not yet published its report.
If such a failure is returned as an error, the model often stops or, worse, guesses a value.
Every answer of the server therefore uses the same envelope.
Besides the data, it states whether the data is available, empty or unavailable; for unavailable data it gives the reason, whether a retry may help and after how long, and which other tools may serve the same need.
It also states the units of every number and the provenance of the data: the source, the time of the data and whether the answer came from the cache.
With this contract the model can choose its next step itself, and the final answer can cite where each number came from.

\subsection{Control Over Actions}\label{sec:action-control}
Most tools of the data server only read data, and each tool declares whether it reads or writes.
The harness runs reading tools at once.
Tools that change something, such as saving a watch list or writing a file, are paused until the user confirms them, and a harness running without a user refuses them.
The user can allow a tool once, for the session or permanently, and a permanent permission is tied to the exact server, so a different server with the same tool name is asked again.

\section{Data Flow and Processing}\label{sec:dataflow}
\Cref{fig:mcp-lifecycle} follows one tool call from the harness to the data providers and back.
\begin{figure}[htb]
  \centering
  \includegraphics[width=0.82\textwidth]{mcp-request-lifecycle}
  \caption{Life cycle of one tool call through the opentrading-mcp server}\label{fig:mcp-lifecycle}
  \fignote{Source: the group, drawn from opentrading-mcp \texttt{de956f8}; diagram source \texttt{figures/src/mcp-request-lifecycle.tex}.}
\end{figure}
The call first passes an HTTP layer that protects the server.
The request is logged; it is refused at once if too many calls are already waiting, so that an overloaded server answers quickly instead of timing out; and it must carry a valid access token, stay within a daily quota and respect a rate limit per client.
A refused call returns a clear reason and, where relevant, the time after which a retry is useful.

The call then enters the tool runtime.
The runtime first looks for a stored answer to the same question.
How long an answer may be reused depends on the market: prices are kept only briefly during a trading session and much longer outside it, and an answer that a provider could not give is remembered for a short time, so that the server does not ask the same failing provider again and again.
If an identical call is already running, the new call waits for its result instead of starting a second request to the provider.
Otherwise the work is given to a pool of workers; long computations, such as backtests, run in a separate pool so that they cannot block short lookups.
Each provider is wrapped by a circuit breaker: after repeated failures the server stops calling that provider for a while and returns an unavailable answer at once.
Before the result is returned, the runtime applies the cleaning rules of \cref{sec:market-data} and builds the answer envelope of \cref{sec:answer-contract}, and the answer is stored for the next identical call.

Some data does not wait for a request.
Background tasks keep live order-book and trade data from the stock exchange during trading hours only; they pause at the lunch break and do not run on public holidays.
In the benchmark, the same tools are called, but their answers come from the frozen recording of \cref{sec:evaluation-data} instead of the providers.

\section{Deployment Strategy}\label{sec:deployment}
The two products are deployed in different ways because they are used in different ways.

OpenTrading runs on the user's own computer.
Each release is built into one program file per operating system, with its runtime included, so that the user does not install any other software.
The files are published as releases of the repository and also through the npm package registry, and the user installs them with a one-line script; the program can update itself to the latest release.
A release is built and checked by automated workflows on real Windows, macOS and Linux machines before it is published.

The data server is packaged as a container image.
For development it runs on the developer's computer and talks to the harness directly through standard input and output, without a network or a token.
For the group it runs as a web service in the cloud, together with a PostgreSQL database for stored statistics and a Redis store for the cache and the shared state of the background tasks.
A local setup with the same three services is described in one compose file, so the same image runs in both places.
Two health addresses support operation: one answers only whether the process is alive, and the other checks the database, the cache and the background worker.
Credentials for the providers and the access tokens are given to the service as environment settings and are never stored in the repository.
\todo{Describe the service as public only after the supervisor confirms the data licences.}

\section{Scalability and Maintenance}\label{sec:scalability}
\subsection{Scalability}\label{sec:scaling}
The data server is designed to stay responsive when many calls arrive together, rather than to process an unlimited number of them.
Three mechanisms serve this goal.
Connections to providers are shared and reused instead of opened for every call.
The number of calls that may wait or run at the same time is bounded; above the bound new calls are refused at once with a retry time.
Identical calls are merged, and answers are reused for as long as the market allows, which avoids asking a provider twice for the same data.
The group did not run a new load test for this version, so no throughput figure is claimed here.

\subsection{Maintenance}\label{sec:maintenance}
Both products are maintained through the same process: every change is a pull request that must pass automated tests before it is merged.
At the versions described in this report, OpenTrading passed \num{2253} tests and opentrading-mcp passed \num{1338} tests with 71.69\% line coverage.
%% Evidence: E11, E13.
OpenTrading follows numbered releases with a change log, and the data server records its changes in a dated change log.
The data server also tracks its dependencies: an audit in November 2026 found that the library used for Vietnamese market data is published only on its vendor's own package index, so a failure of that index would stop every build.
This risk, together with open work on the handling of provider sessions and health checks, is listed as future work in \cref{sec:future}.
